\documentclass[12pt,a4paper]{article}
\usepackage{amsmath,amssymb,amsthm,hyperref,geometry}
\usepackage{enumitem,booktabs,array}
\usepackage[T1]{fontenc}
\usepackage[utf8]{inputenc}
\geometry{margin=2.5cm}

% -------------------------------------------------------
% Theorem environments
% -------------------------------------------------------
\newtheorem{theorem}{Theorem}[section]
\newtheorem{lemma}[theorem]{Lemma}
\newtheorem{corollary}[theorem]{Corollary}
\newtheorem{proposition}[theorem]{Proposition}
\newtheorem{conjecture}[theorem]{Conjecture}
\newtheorem{definition}[theorem]{Definition}

\theoremstyle{remark}
\newtheorem{remark}[theorem]{Remark}
\newtheorem{example}[theorem]{Example}

% -------------------------------------------------------
% Metadata
% -------------------------------------------------------
\title{\textbf{The $abc$ Conjecture}\\[0.5em]
\large Radical, Height, Mochizuki's IUT Theory, and Consequences}
\author{Michael Fuhrmann}
\date{2026-03-11}

\hypersetup{
  pdftitle={The abc Conjecture},
  pdfauthor={Michael Fuhrmann},
  colorlinks=true,
  linkcolor=blue,
  citecolor=blue
}

\begin{document}
\maketitle
\begin{abstract}
The $abc$ conjecture, proposed independently by Oesterlé and Masser in 1985, is one of
the most far-reaching open problems in number theory.  It posits a profound inequality
relating the sizes of three coprime integers $a,b,c$ with $a+b=c$ to the product of
their distinct prime factors—the \emph{radical} $\operatorname{rad}(abc)$.  The conjecture
implies, as special cases, Fermat's Last Theorem for large exponents, Roth's theorem,
Szpiro's conjecture on elliptic curves, and many classical Diophantine results.

In 2012 Shinichi Mochizuki announced a proof via his \emph{Inter-universal Teichmüller
theory} (IUT); as of 2026 the mathematical community remains divided, with the
Scholze--Stix objection (2018) unresolved.  This paper surveys the conjecture's
formulation, its most important consequences, the strategy of Mochizuki's approach, and
the current state of the controversy.
\end{abstract}

\tableofcontents
\newpage

% ================================================================
\section{Formulation of the Conjecture}
% ================================================================

\subsection{The Radical}

\begin{definition}[Radical]
  For a positive integer $n$, the \emph{radical} of $n$ is
  \[
    \operatorname{rad}(n) \;:=\; \prod_{\substack{p\;\text{prime} \\ p \mid n}} p.
  \]
  In other words, $\operatorname{rad}(n)$ is the largest squarefree divisor of $n$.
\end{definition}

\begin{example}
  $\operatorname{rad}(12) = \operatorname{rad}(2^2\cdot 3) = 2\cdot 3 = 6$;
  $\quad \operatorname{rad}(360) = \operatorname{rad}(2^3\cdot 3^2\cdot 5) = 30$.
\end{example}

\subsection{Height Function}

The \emph{logarithmic height} (or \emph{quality}) of a triple $(a,b,c)$ with $a+b=c$,
$\gcd(a,b)=1$, $a,b,c>0$ is defined as
\begin{equation}\label{eq:quality}
  q(a,b,c) \;:=\; \frac{\log\max(a,b,c)}{\log\operatorname{rad}(abc)}
             \;=\; \frac{\log c}{\log\operatorname{rad}(abc)}.
\end{equation}
The $abc$ conjecture essentially asserts that the quality $q(a,b,c)$ cannot be much
larger than $1$.

\subsection{The $abc$ Conjecture}

\begin{conjecture}[$abc$ Conjecture, Oesterlé--Masser 1985~\cite{Oesterle1988,Masser1985}]
  For every $\varepsilon > 0$ there exists a constant $K_\varepsilon > 0$ such that, for
  all triples $(a,b,c)$ of coprime positive integers with $a + b = c$,
  \begin{equation}\label{eq:abc}
    c \;\leq\; K_\varepsilon \cdot \operatorname{rad}(abc)^{1+\varepsilon}.
  \end{equation}
\end{conjecture}

Equivalently, only finitely many triples have quality $q(a,b,c) > 1+\varepsilon$.

\begin{remark}
  The $\varepsilon$ is necessary: the record triple $a=1$, $b=3^{10}\cdot109$,
  $c=23^5$ has quality $q \approx 1.6299 > 1$, so $c = \operatorname{rad}(abc)^{1.6299}
  \gg \operatorname{rad}(abc)$, showing the bound $c \leq K \cdot \operatorname{rad}(abc)$
  (exponent $1$ without $\varepsilon$) cannot hold universally.  Triples with quality
  $q \to 1$ from below exist in abundance (e.g.\ $a=1$, $b=2^n-1$, $c=2^n$ when
  $2^n-1$ is a Mersenne prime) but do \emph{not} violate the $\varepsilon=0$ bound;
  it is the record triples with $q > 1$ that force $\varepsilon > 0$.  The largest
  known quality exceeding $1$ is $q \approx 1.6299$.
\end{remark}

\subsection{Equivalent Formulations}

An equivalent form uses the height $h(a,b,c) = \log\max(|a|,|b|,|c|)$: the $abc$
conjecture states
\[
  h(a,b,c) \;\leq\; (1+\varepsilon)\sum_{p\mid abc} \log p + O_\varepsilon(1)
\]
for all coprime $(a,b,c)$ with $a+b+c=0$ (in the projective formulation with signed integers).

A weak form asserts the existence of $K$ (depending on $\varepsilon$) such that for
\emph{all} coprime triples:
\[
  \max(|a|,|b|,|c|) \;\leq\; K \cdot \operatorname{rad}(abc)^2.
\]
This weak $abc$ conjecture (exponent $2$ in place of $1+\varepsilon$) follows from the
theory of linear forms in logarithms (Baker's theorem), making it unconditionally true.

% ================================================================
\section{High-Quality $abc$ Triples}
% ================================================================

Systematic computer searches (Reyssat, Nitaj, \cite{Nitaj1996}) have catalogued triples
with large quality.  The ten highest-quality known triples all have $q > 1.4$:

\begin{table}[h]
\centering
\begin{tabular}{lll}
\toprule
$a$ & $b$ & $c = a+b$ \quad (quality) \\
\midrule
$1$ & $2$ & $3$ \quad ($q=1$, trivial) \\
$1$ & $8$ & $9 = 3^2$ \quad ($q = \log9/\log6 \approx 1.226$) \\
$2^5$ & $7^2$ & $3^4$ \quad ($q\approx1.292$, all small) \\
$1$ & $3^{10}\!\cdot\!109$ & $23^5$ \quad ($q\approx1.630$, current record) \\
\bottomrule
\end{tabular}
\caption{Selected high-quality $abc$ triples}
\end{table}

The Goldfeld conjecture, a refinement, predicts $N(Q,X) \sim C_\varepsilon X^{2-2/Q}$ for the
number of triples with $c\leq X$ and quality $>Q$.

% ================================================================
\section{Consequences of the $abc$ Conjecture}
% ================================================================

\subsection{Fermat's Last Theorem for Large Exponents}

\begin{theorem}[$abc$ $\Rightarrow$ FLT for $n \geq 6$~\cite{Granville1990}]
  Assume the $abc$ conjecture.  Then for all $n \geq 6$, the equation $x^n + y^n = z^n$
  has no solution in positive integers.
\end{theorem}

\begin{proof}[Proof]
  Suppose $x^n + y^n = z^n$ with $\gcd(x,y)=1$.  Set $a=x^n$, $b=y^n$, $c=z^n$.
  Then $\operatorname{rad}(abc) \leq xyz \leq z^3$.  The $abc$ conjecture (with
  $\varepsilon=1$) gives
  \[
    c = z^n \;\leq\; K_1 \cdot (z^3)^2 = K_1 z^6,
  \]
  hence $z^{n-6}\leq K_1$, which bounds $z$ (and hence $x,y$) for $n>6$.  Refining with
  $\varepsilon < 1$ extends the argument to $n\geq 6$.
\end{proof}

\begin{remark}
  Wiles's proof of FLT (1995) is valid for all $n\geq 3$ and is unconditional.
  The $abc$-based proof is much simpler but requires the unproved $abc$ conjecture and
  only covers $n\geq 6$ (or with more work, $n\geq 3$).
\end{remark}

\subsection{Roth's Theorem and Thue--Mahler Equations}

\begin{corollary}[$abc$ $\Rightarrow$ Roth's theorem]
  The $abc$ conjecture implies: for any algebraic number $\alpha$ of degree $d\geq 2$ and
  any $\varepsilon>0$, there are only finitely many rationals $p/q$ with
  $|\alpha - p/q| < q^{-2-\varepsilon}$.
\end{corollary}

More precisely, $abc$ implies the full strength of Baker's theory of linear forms in
logarithms (and more), making all effective Diophantine approximation results uniform.

\subsection{Szpiro's Conjecture for Elliptic Curves}

Let $E/\mathbb{Q}$ be an elliptic curve with minimal discriminant $\Delta_{\min}$ and
conductor $N$.

\begin{conjecture}[Szpiro's Conjecture~\cite{Szpiro1981}]
  For every $\varepsilon>0$ there exists $C_\varepsilon$ such that
  \begin{equation}\label{eq:szpiro}
    |\Delta_{\min}| \;\leq\; C_\varepsilon \cdot N^{6+\varepsilon}.
  \end{equation}
\end{conjecture}

\begin{theorem}[Frey--Masser--Oesterlé; see~\cite{Silverman1994}]
  The $abc$ conjecture is equivalent to Szpiro's conjecture \eqref{eq:szpiro} for
  semistable elliptic curves.
\end{theorem}

The connection is the \emph{Frey curve}: given an $abc$ triple, the elliptic curve
$y^2 = x(x-a^n)(x+b^n)$ has discriminant related to $(abc)^{2n}$ and conductor
essentially $\operatorname{rad}(abc)^2$.

\subsection{Further Consequences}

The $abc$ conjecture implies or simplifies proofs of:
\begin{itemize}[nosep]
  \item \textbf{Mordell--Faltings:} The Faltings-type bound $h(P)\leq c(E)$
    for heights of rational points.
  \item \textbf{Wieferich primes:} Only finitely many Wieferich primes
    (primes $p$ with $2^{p-1}\equiv1\pmod{p^2}$) -- conditional on $abc$
    (Granville--Langevin 1985).
  \item \textbf{Hall's conjecture:} $|x^3 - y^2| \gg x^{1/2 - \varepsilon}$
    for distinct integers, implied by $abc$.
  \item \textbf{Catalan's conjecture:} Only $2^3$ and $3^2$ are consecutive perfect powers
    (proved unconditionally by Mihailescu 2002, but simpler under $abc$).
  \item \textbf{Generalised Fermat:} $x^p + y^q = z^r$ has finitely many coprime
    solutions for $1/p+1/q+1/r < 1$.
\end{itemize}

% ================================================================
\section{Mochizuki's Inter-Universal Teichmüller Theory (IUT)}
% ================================================================

\subsection{Overview}

In four preprints posted in August 2012~\cite{Mochizuki2012a,Mochizuki2012b,
Mochizuki2012c,Mochizuki2012d}, Shinichi Mochizuki (RIMS, Kyoto) announced a proof of
the $abc$ conjecture via a radically new framework he calls
\emph{Inter-Universal Teichmüller Theory} (IUT, also called \emph{arithmetic deformation
theory}).

\subsection{Key Ideas of IUT}

IUT operates with objects called \emph{Hodge theatres}, which are elaborate combinatorial
and categorical structures built from data associated to an elliptic curve over a number
field.  The central construction involves:

\begin{enumerate}[nosep]
  \item \textbf{Theatres and links:}  A \emph{$\Theta^{\pm\text{ell}}$NF-Hodge theatre}
    packages the arithmetic of the elliptic curve and an auxiliary prime $l$ into a
    categorical object.  \emph{Étale theta functions} encode the $p$-adic geometry.
  \item \textbf{Multiradiality:}  Mochizuki introduces a distinction between
    \emph{mono-anabelian} and \emph{poly-anabelian} structures.  The key principle is
    that certain data can be reconstructed from the absolute Galois group alone
    (\emph{mono-anabelic reconstruction}).
  \item \textbf{Log-theta lattice:}  Successive applications of the \emph{log-link}
    (which passes from multiplicative to additive structures via $p$-adic logarithms)
    and the \emph{theta-link} (relating different theatres) form a lattice.
  \item \textbf{Kummer theory and indeterminacies:}  Three types of indeterminacy
    (labeled $\mathrm{Ind}1$, $\mathrm{Ind}2$, $\mathrm{Ind}3$) arise from the
    non-commutativity of the log- and theta-links.  Controlling these indeterminacies
    via explicit bounds yields the $abc$ inequality.
\end{enumerate}

The strategy of the proof is to show:
\[
  -\log|\Delta_E| \;\leq\; (1+\varepsilon)\log\operatorname{rad}(\Delta_E) + O(1),
\]
which translates (via Szpiro's equivalence) into the $abc$ conjecture.

\subsection{The Scholze--Stix Objection (2018)}

In March 2018, Peter Scholze and Jakob Stix visited Mochizuki in Kyoto and subsequently
wrote a critical report~\cite{ScholzeStix2018} identifying what they believe is a
fundamental error in IUT.

Their central objection concerns Mochizuki's claim that certain identifications of
ring structures (which he refers to as the ``theta-link'') can be made compatible with the
ring structure on both sides.  Scholze--Stix argue that after the theta-link, the relevant
inequality reduces to a trivial identity $\log(q) \leq \log(q)$ and therefore carries no
content.  In their view, the indeterminacies are too coarse to establish a non-trivial
bound.

Mochizuki responded with detailed rebuttals~\cite{Mochizuki2018rebuttal}, maintaining
that Scholze--Stix misidentified the roles of the distinct ``copies'' of the ring.

\subsection{Current Situation (2026)}

The Japanese journal \textit{Publications of RIMS} published Mochizuki's four IUT papers
in March 2021 after an internal review process.  However, this has not resolved the
controversy:
\begin{itemize}[nosep]
  \item The Scholze--Stix objection is still considered valid by many leading number
    theorists (Scholze, Stix, Conrad, de Jong, \ldots).
  \item A small group of researchers (Fesenko, Yamashita, Minamide, \ldots) have studied
    IUT in depth and affirm its correctness.
  \item No third-party, independent verification by a broad consensus has emerged.
  \item A condensed/simplified account of IUT remains unavailable.
\end{itemize}

% BUG-B8-P34-STATUS behoben: abc-Vermutung explizit als OFFEN klassifiziert
\begin{remark}[\textbf{Official status of the abc conjecture — March 2026}]
\textbf{The $abc$ conjecture is \emph{OPEN} as of March 2026.}
\begin{itemize}[nosep]
  \item Mochizuki's IUT papers were published in PRIMS in 2021, but publication in a
    journal does \emph{not} constitute proof accepted by the community.
  \item The Scholze--Stix objection (2018) is considered valid by the
    \emph{majority} of the international algebraic geometry and number theory community.
    Key experts who have publicly endorsed the objection or remain unconvinced of IUT
    include Scholze, Stix, Conrad, de Jong, Venkatesh, and others.
  \item No independent, broadly accepted verification of IUT has appeared.
  \item \textbf{Consequence:} Any result that ``follows from the abc conjecture'' is
    still conditional.  The abc conjecture should be cited as an open problem, not as
    an established theorem.
\end{itemize}
\end{remark}

The mathematical community is effectively in an impasse, and the $abc$ conjecture is
\textbf{not proved}.

% ================================================================
\section{Alternative Approaches}
% ================================================================

\subsection{$p$-Adic and Arakelov Geometry}

Several researchers have explored whether elements of IUT can be replaced by more
conventional $p$-adic Hodge theory or Arakelov intersection theory.  Tan~\cite{Tan2022}
and others have sketched partial reformulations; none have so far succeeded in bypassing
the theta-link.

\subsection{Function Field Analogue}

Over function fields $k(C)$ (where $C$ is a curve over a finite field $\mathbb{F}_q$),
the $abc$ conjecture is a theorem, following from the Mason--Stothers theorem
(Mason 1983~\cite{Mason1984}):

\begin{theorem}[Mason--Stothers Theorem]
  Let $k$ be a field of characteristic $0$.  If $a,b,c\in k[t]$ are coprime polynomials
  with $a+b=c$ and $\max(\deg a,\deg b,\deg c)>0$, then
  \[
    \max(\deg a,\deg b,\deg c) \;\leq\; \deg\operatorname{rad}(abc) - 1.
  \]
\end{theorem}

This is significantly stronger than the number field $abc$ conjecture (the exponent is
exactly $1$, not $1+\varepsilon$) and is proved by a simple Wronskian argument.  The
analogy suggests that an arithmetic analogue of the Wronskian might be the key to
the number field case.

% ================================================================
\section{The $abc$ Conjecture and the Langlands Programme}
% ================================================================

The $abc$ conjecture can be formulated in the language of heights on the projective line
$\mathbb{P}^1$ minus three points, relating to the \emph{Hyperbolic Mordell} or
\emph{Shafarevich--Faltings} setting.  Connections include:

\begin{itemize}[nosep]
  \item \textbf{Vojta's conjecture:}  A generalisation of $abc$ to rational points on
    higher-dimensional varieties; implies Faltings' theorem and more.
  \item \textbf{$S$-unit equation:}  Over number fields, $\{a+b=c : a,b,c\in\mathcal{O}_S^*\}$
    is finite (Thue--Mahler--$S$-unit theorem); $abc$ gives quantitative bounds.
  \item \textbf{Weil height machine:}  The height $h(a,b,c)$ is a Weil height on
    $\mathbb{P}^1$; Vojta's conjecture in dimension $1$ is exactly $abc$.
\end{itemize}

% ================================================================
\section{Current State of Proof}
% ================================================================

\begin{itemize}[nosep]
  \item \textbf{Proved:} Weak $abc$ (exponent $2$, via Baker), Mason--Stothers (function fields),
    Szpiro's for CM curves (Goldfeld--Szpiro), Fermat $n\geq 3$ (Wiles, unconditional).
  \item \textbf{Conditional on $abc$:} FLT for $n\geq 6$ (much simpler), Wieferich primes
    (finitely many), Hall's conjecture, generalised Catalan.
  \item \textbf{Mochizuki's IUT (2012):} Claimed proof; published in PRIMS 2021; consensus
    not achieved; Scholze--Stix objection unresolved.
  \item \textbf{Open:} The $abc$ conjecture is not considered proved.
\end{itemize}

% ================================================================
\section{Conclusion}
% ================================================================

The $abc$ conjecture occupies a singular position in modern number theory: a single
inequality with enormous reach, connecting additive and multiplicative structures of
integers in a way that is both elementary to state and extraordinarily deep.  Its proof
would reorganise wide swaths of Diophantine geometry.  Whether Mochizuki's IUT
eventually achieves consensus acceptance or a new approach emerges remains to be seen.

\begin{thebibliography}{99}

\bibitem{Oesterle1988}
J.~Oesterlé,
\textit{Nouvelles approches du théorème de Fermat},
Séminaire Bourbaki, Exposé \textbf{694} (1988), 165--186.

\bibitem{Masser1985}
D.~W.~Masser,
\textit{Open problems},
in \textit{Proc.\ Sympos.\ Analytic Number Theory} (W.~W.~L.~Chen, ed.),
Imperial College, London, 1985.

\bibitem{Granville1990}
A.~Granville,
\textit{The Beal conjecture and the $abc$ conjecture},
unpublished note, 1990.
See also: A.~Granville and T.~J.~Tucker,
\textit{It's as easy as $abc$},
Notices AMS \textbf{49} (2002), 1224--1231.

\bibitem{Szpiro1981}
L.~Szpiro,
\textit{Séminaire sur les pinceaux de courbes de genre au moins deux},
Astérisque \textbf{86} (1981).

\bibitem{Silverman1994}
J.~H.~Silverman,
\textit{Advanced Topics in the Arithmetic of Elliptic Curves},
Springer, GTM \textbf{151}, 1994.

\bibitem{Nitaj1996}
A.~Nitaj,
\textit{La conjecture $abc$},
Enseign.\ Math.\ \textbf{42} (1996), 3--24.

\bibitem{Mason1984}
R.~C.~Mason,
\textit{Diophantine Equations over Function Fields},
LMS Lecture Note Series \textbf{96}, Cambridge University Press, 1984.

\bibitem{Mochizuki2012a}
S.~Mochizuki,
\textit{Inter-Universal Teichmüller Theory I: Construction of Hodge Theaters},
preprint (2012), RIMS Kyoto.

\bibitem{Mochizuki2012b}
S.~Mochizuki,
\textit{Inter-Universal Teichmüller Theory II: Hodge-Arakelov-Theoretic Evaluation},
preprint (2012), RIMS Kyoto.

\bibitem{Mochizuki2012c}
S.~Mochizuki,
\textit{Inter-Universal Teichmüller Theory III: Canonical Splittings of the Log-Theta-Lattice},
preprint (2012), RIMS Kyoto.

\bibitem{Mochizuki2012d}
S.~Mochizuki,
\textit{Inter-Universal Teichmüller Theory IV: Log-Volume Computations and Set-Theoretic Foundations},
preprint (2012), RIMS Kyoto.

\bibitem{ScholzeStix2018}
P.~Scholze and J.~Stix,
\textit{Why abc is still a conjecture},
report (2018), available at \url{https://www.math.uni-bonn.de/people/scholze/WhyABCisStillaConjecture.pdf}.

\bibitem{Mochizuki2018rebuttal}
S.~Mochizuki,
\textit{Comments on the manuscript by Scholze--Stix},
report (2018), available at RIMS Kyoto website.

\bibitem{Tan2022}
F.~Tan,
\textit{A remark on Mochizuki's Corollary 3.12},
preprint, 2022.

\bibitem{Vojta1987}
P.~Vojta,
\textit{Diophantine Approximations and Value Distribution Theory},
Lecture Notes in Mathematics \textbf{1239}, Springer, 1987.

\end{thebibliography}

\end{document}
